% titel: Blatt 10, A6: Lewis-Formeln und pH der Chloressigsäuren
% Teil II: der durchgerechnete Weg.

\subsection*{a) Lewis-Strukturen (Skizze, repräsentativ)}
Alle vier sind Carbonsäuren mit der Carboxyl-Gruppe
$R$--C(=O)--OH und zunehmender Chlorierung am $\alpha$-C:
$\mathrm{CH_3COOH},\ \mathrm{CH_2ClCOOH},\ \mathrm{CHCl_2COOH},\
\mathrm{CCl_3COOH}$. Die Struktur zeigt um das carbonyl C eine C=O-Doppel-
bindung, eine O--H und die R-Gruppe; Chlor ersetzt schrittweise die
Wasserstoff am $\alpha$-C.

\subsection*{b) pH-Werte bei $c_0=0.01\ M$}
Näherung $pH\approx\frac12(pK_s+2)$ (da $-\lg 0.01=2$):
\begin{align*}
  \mathrm{CH_3COOH{:}\;} & pK_s=4.75
    \Rightarrow pH \approx \frac{4.75+2}{2} \approx 3.38,\\
  \mathrm{CH_2ClCOOH{:}\;} & pK_s=2.85
    \Rightarrow pH \approx \frac{2.85+2}{2} \approx 2.42,\\
  \mathrm{CHCl_2COOH{:}\;} & pK_s=1.48
    \Rightarrow pH \approx \frac{1.48+2}{2} \approx 1.74,\\
  \mathrm{CCl_3COOH{:}\;} & pK_s=0.52
    \Rightarrow pH \approx \frac{0.52+2}{2} \approx 1.26 .
\end{align*}

\begin{ergebnis}
pH(0.01\,M): Essigsäure $\approx3.4$; Mono-Chlor $\approx2.4$; Di-Chlor
$\approx1.7$; Tri-Chlor $\approx1.3$. Jede Cl-Gruppe macht die Lösung
sauer (stärker sauer).
\end{ergebnis}

\begin{kontrolle}
Die Reihenfolge folgt den $pK_s$ (kleiner $pK_s\Rightarrow$ kleinerer pH).
Plausibilität: bei $c_0{=}0.01$ M wäre eine starke Säure pH$=$2; alle Werte
$\geq1.26$ liegen über (schwach bis mäßig stark), Trichlor nahe 2 ist fast
stark. Kein Widerspruch.
\end{kontrolle}

\begin{merksatz}
$pH\approx\frac12(pK_s-\lg c_0)$. Elektronenziehende Cl-Atome am
$\alpha$-C stabilisieren das Anion und stärken die Säure (niedrigerer
pK$_s$).
\end{merksatz}
